\documentclass{article}

%maths
\usepackage{mathtools}
\usepackage{amsmath}
\usepackage{amssymb}
\usepackage{amsfonts}

%tikzpicture
\usepackage{tikz}
\usepackage{scalerel}
\usepackage{pict2e}
\usepackage{tkz-euclide}
\usetikzlibrary{calc}
\usetikzlibrary{patterns,arrows.meta}
\usetikzlibrary{shadows}

%pgfplots
\usepackage{pgfplots}
\pgfplotsset{compat=newest}
\usepgfplotslibrary{statistics}
\usepgfplotslibrary{fillbetween}

%colours
\usepackage{xcolor}

\usepackage{nicefrac}

\begin{document}

\section{Polynomial Functions}
A \textbf{Polynomial Function} is any function of the form:\\
\begin{center}
	$a_0x^n +a_1x^{n-1}  +a_2x^{n-2} +...+ a_{n-1}x+ a_n$ 
\end{center}

\subsection{Zeros of a Function}
\textbf{Zeros of a Function} is when you make a function equal zero so you can find the roots which equal zero. This tells you the x-intercepts of the graph. \\
\\
\end{document}
